\documentclass[a4paper,10pt,twocolumn]{ltjarticle}
\usepackage{kaitoushu}

% グラフ用の共通スタイル定義
\pgfplotsset{
  myaxis/.style={
    axis lines=middle,
    xlabel={$x$}, ylabel={$y$},
    every axis x label/.style={at={(ticklabel* cs:1)},anchor=west},
    every axis y label/.style={at={(ticklabel* cs:1)},anchor=south},
    tick label style={font=\scriptsize},
    label style={font=\small},
    width=5.5cm, height=5cm,
    grid=none,
    samples=100,
  }
}

\begin{document}

\problemrange{4章：指数関数と対数関数 — Check (p.51) 対数（問240--247）}



\mondai{240}{次の等式を満たす $x$ の値を求めよ．}
\kotae{
対数の定義 $\log_{a}x=c$ $\Leftrightarrow$ $x=a^{c}$ を使う．

(1) $\log_{4}x=3$ → $x=4^{3}=64$

(2) $\log_{x}9=2$ → $x^{2}=9$，$x>0,\,x\neq1$ より $x=3$

(3) $\log_{0.1}x=-2$ → $x=0.1^{-2}=100$

(4) $\log_{x}\dfrac{1}{5}=\dfrac{1}{3}$ → $x^{1/3}=\dfrac{1}{5}$，$x=\left(\dfrac{1}{5}\right)^{3}=\dfrac{1}{125}$
}

\mondai{241}{次の式を計算せよ．}
\kotae{
(1) $\log_{6}18+\log_{6}72=\log_{6}(18\times72)=\log_{6}1296=\log_{6}6^{4}=4$

(2) $\log_{2}432-3\log_{2}6=\log_{2}432-\log_{2}6^{3}=\log_{2}\dfrac{432}{216}=\log_{2}2=1$

(3) $3\log_{2}\sqrt{12}+\log_{2}\dfrac{2}{3\sqrt{3}}$

$=\log_{2}(\sqrt{12})^{3}+\log_{2}\dfrac{2}{3\sqrt{3}}=\log_{2}\!\left(12\sqrt{12}\cdot\dfrac{2}{3\sqrt{3}}\right)$

$=\log_{2}\!\left(\dfrac{2\cdot12\sqrt{12}}{3\sqrt{3}}\right)=\log_{2}\!\left(\dfrac{24\cdot2\sqrt{3}}{3\sqrt{3}}\right)=\log_{2}16=4$

(4) $\log_{5}\sqrt[3]{15}+\log_{5}\sqrt{5}-\log_{5}\sqrt[3]{3}$

$=\log_{5}\dfrac{\sqrt[3]{15}\cdot\sqrt{5}}{\sqrt[3]{3}}=\log_{5}\left(\sqrt[3]{5}\cdot5^{1/2}\right)=\log_{5}\,5^{1/3+1/2}=\log_{5}5^{5/6}=\dfrac{5}{6}$

(5) $\left(\log_{3}\dfrac{1}{4}\right)\cdot(\log_{2}3)=(-2\log_{3}2)\cdot\log_{2}3=-2\cdot\dfrac{\log 2}{\log 3}\cdot\dfrac{\log 3}{\log 2}=-2$

(6) $(\log_{\sqrt{2}}3)\cdot(\log_{9}8)=\dfrac{\log 3}{\log\sqrt{2}}\cdot\dfrac{\log 8}{\log 9}=\dfrac{\log 3}{\frac{1}{2}\log 2}\cdot\dfrac{3\log 2}{2\log 3}=\dfrac{2\log 3}{\log 2}\cdot\dfrac{3\log 2}{2\log 3}=3$
}

\mondai{242}{次の関数のグラフは $y=\log_{2}x$ のグラフをどのように移動したものか．}
\kotae{
(1) $y=\log_{2}(x-1)+2$

$x$ を $x-1$ に → $x$軸方向に$1$平行移動，$+2$ → $y$軸方向に$2$平行移動

(2) $y=-\log_{2}(-x)$

$-x$ → $y$軸対称，$-$ → $x$軸対称．合わせて原点に関して対称移動．
}

\mondai{243}{次の方程式を解け．}
\kotae{
(1) $\log_{4}(x-1)+\log_{4}(x-7)=2$

真数条件：$x>7$

$(x-1)(x-7)=4^{2}=16$，$x^{2}-8x+7=16$，$x^{2}-8x-9=0$

$(x-9)(x+1)=0$．$x>7$ より $x=9$

(2) $\log_{5}(x^{2}+6x-7)-\log_{5}(x-1)=2$

真数条件：$x>1$

$\log_{5}\dfrac{x^{2}+6x-7}{x-1}=2$

$\dfrac{(x+7)(x-1)}{x-1}=x+7=25$ → $x=18$\quad（$>1$で適）

(3) $2\log_{6}(x+2)-\log_{6}x=2\log_{6}3$

$\log_{6}(x+2)^{2}=\log_{6}(9x)$

$(x+2)^{2}=9x$，$x^{2}+4x+4=9x$，$x^{2}-5x+4=0$

$(x-1)(x-4)=0$．真数条件 $x>0$ よりどちらも適．$x=1,\,4$
}

\mondai{244}{次の不等式を解け．}
\kotae{
(1) $\log_{4}(3x-1)<\dfrac{3}{2}$

真数条件：$x>\dfrac{1}{3}$

$3x-1<4^{3/2}=8$，$3x<9$，$x<3$

$\therefore \dfrac{1}{3}<x<3$

(2) $\log_{1/3}(x+1)<-2$

真数条件：$x>-1$

底 $\dfrac{1}{3}<1$ なので不等号反転：$x+1>\left(\dfrac{1}{3}\right)^{-2}=9$

$x>8$
}

\mondai{245}{$\log_{10}2=a$，$\log_{10}3=b$ のとき，次の式を $a,b$ で表せ．}
\kotae{
基本：$\log_{10}5=\log_{10}\dfrac{10}{2}=1-a$

(1) $\log_{10}12=\log_{10}(2^{2}\cdot3)=2a+b$

(2) $\log_{10}0.3=\log_{10}\dfrac{3}{10}=b-1$

(3) $\log_{5}6=\dfrac{\log_{10}6}{\log_{10}5}=\dfrac{a+b}{1-a}$

(4) $\log_{8}15=\dfrac{\log_{10}15}{\log_{10}8}=\dfrac{\log_{10}3+\log_{10}5}{3\log_{10}2}=\dfrac{b+(1-a)}{3a}=\dfrac{1-a+b}{3a}$
}

\mondai{246}{次の不等式を満たす最大の整数 $n$ を求めよ．ただし，$\log_{10}2=0.3010$ とする．$0.1^{n}\geqq\dfrac{1}{2^{100}}$}
\kotae{
$10^{-n}\geqq2^{-100}$

両辺の常用対数：$-n\geqq-100\log_{10}2=-30.10$

$n\leqq30.10$ → $n=30$
}

\mondai{247}{ある国の今年の人口は昨年より3\%減少した．来年以降も同じ割合で減少すると，今年の人口の半分以下となるのは何年後か．ただし，$\log_{10}2=0.3010$，$\log_{10}9.7=0.9868$ とする．}
\kotae{
$n$ 年後の人口：$(0.97)^{n}$ 倍．

$(0.97)^{n}\leqq0.5$ の両辺の常用対数：

$n\log_{10}0.97\leqq\log_{10}0.5$

$\log_{10}0.97=\log_{10}\dfrac{9.7}{10}=0.9868-1=-0.0132$

$\log_{10}0.5=\log_{10}\dfrac{1}{2}=-0.3010$

$-0.0132\,n\leqq-0.3010$（両辺を負で割るので不等号反転）

$n\geqq\dfrac{0.3010}{0.0132}\approx22.8$

$\therefore$ 23年後
}

\end{document}
